% figures/fig5_geometria_2d_angular.tex
\begin{tikzpicture}[
    >=Stealth,
    sensor/.style={
        circle, draw=black!85, fill=black!12, line width=1pt,
        minimum size=0.52cm, inner sep=0pt
    },
    source/.style={
        circle, draw=black!85, fill=orange!80, line width=1.1pt,
        minimum size=0.55cm, inner sep=0pt
    },
    axis/.style={
        ->, line width=0.85pt, draw=black!65
    }
]

    % =========================================================================
    % 1. EJES CARTESIANOS (Origen en x = 0.5)
    % =========================================================================
    \coordinate (Orig) at (0.5, 0.0);

    % Ejes coordenados alargados
    \draw[axis] (-5.6, 0.0) -- (5.6, 0.0) node[right=2pt, font=\small\sffamily] {$x$ (Transversal)};
    \draw[axis] (0.5, -0.6) -- (0.5, 5.7) node[above=2pt, font=\small\sffamily] {$y$ (Broadside)};

    % =========================================================================
    % 2. SENSORES RECEPTORES: R1(-d/2, 0) y R2(+d/2, 0)
    % =========================================================================
    % d/2 = 2.0 => R1 en -1.5; R2 en 2.5
    \coordinate (R1) at (-1.5, 0.0);
    \coordinate (R2) at (2.5, 0.0);

    \node[sensor] (NodeR1) at (R1) {};
    \node[below=6pt, font=\small\sffamily] at (R1) {$R_1\left(-\frac{d}{2}, 0\right)$};

    \node[sensor] (NodeR2) at (R2) {};
    \node[below=6pt, font=\small\sffamily] at (R2) {$R_2\left(+\frac{d}{2}, 0\right)$};

    % Cota de línea de base d (inferior con amplio respiro)
    \draw[<->, draw=black!70, line width=0.65pt] (-1.5, -1.30) -- node[below=2pt, font=\footnotesize\sffamily] {Línea de base $d = \norm{\vec{r}_2 - \vec{r}_1}$} (2.5, -1.30);
    \draw[dotted, line width=0.5pt, draw=black!45] (-1.5, 0.0) -- (-1.5, -1.30);
    \draw[dotted, line width=0.5pt, draw=black!45] (2.5, 0.0) -- (2.5, -1.30);

    % =========================================================================
    % 3. FUENTE EMISORA S(x, y) (Ondas hacia abajo, texto limpio arriba)
    % =========================================================================
    \coordinate (S) at (1.7, 4.4);

    % Proyecciones ortogonales hacia los ejes
    \draw[dotted, line width=0.65pt, draw=black!55] (1.7, 4.4) -- (1.7, 0.0) node[below=3pt, font=\small\sffamily] {$x$};
    \draw[dotted, line width=0.65pt, draw=black!55] (1.7, 4.4) -- (0.5, 4.4) node[left=2pt, font=\small\sffamily] {$y$};

    % Frentes de onda concéntricos radiados HACIA ABAJO (hacia los receptores)
    \draw[line width=0.75pt, orange!85!black] ($(S)+(215:0.45)$) arc (215:285:0.45);
    \draw[line width=0.75pt, orange!85!black] ($(S)+(215:0.80)$) arc (215:285:0.80);
    \draw[line width=0.75pt, orange!85!black] ($(S)+(215:1.15)$) arc (215:285:1.15);

    % Nodo Fuente (el rótulo queda arriba, totalmente libre de ondas)
    \node[source] (NodeS) at (S) {};
    \node[above=6pt, font=\small\sffamily\bfseries] at (S) {$S(x, y)$};

    % =========================================================================
    % 4. VECTORES DE PROPAGACIÓN r1, r2, r Y ÁNGULO THETA
    % =========================================================================
    % Vector r1
    \draw[->, line width=1.15pt, draw=UdeAGreen!75!black] (NodeR1) -- node[pos=0.45, above left=1pt, font=\small] {$\vec{r}_1$} (NodeS);

    % Vector r2
    \draw[->, line width=1.15pt, draw=blue!70!black] (NodeR2) -- node[pos=0.55, right=2pt, font=\small] {$\vec{r}_2$} (NodeS);

    % Vector central r
    \draw[->, dashed, line width=0.85pt, draw=black!70] (Orig) -- node[pos=0.65, below right=1pt, font=\small] {$\vec{r}$} (NodeS);

    % Arco del ángulo de incidencia theta respecto al eje vertical y
    \draw[line width=0.8pt, draw=black!85] ($(Orig) + (0.0, 2.2)$) arc (90:74.8:2.2);
    \node[font=\small\sffamily\bfseries] at ($(Orig) + (82.5:2.55)$) {$\theta$};

    % =========================================================================
    % 5. HIPÉRBOLA TDOA
    % =========================================================================
    \draw[thick, orange!90!black, domain=0.4:5.2, smooth, variable=\y]
        plot ({ 0.5 + 0.4845 * sqrt(1 + (\y/1.940)*(\y/1.940)) }, {\y});
    % Rótulo en la parte superior derecha, libre de ondas y de la fuente
    \node[orange!90!black, font=\footnotesize\sffamily, right=3pt] at (1.92, 5.25) {Hipérbola TDOA ($\Delta r = r_1 - r_2$)};

    % =========================================================================
    % 6. RECUADRO ESQUINADO EN EL EXTREMO SUPERIOR IZQUIERDO
    % =========================================================================
    \node[draw=black!75, fill=white, rounded corners=4pt, line width=0.8pt, inner sep=7pt] at (-3.3, 4.2) {
        \footnotesize
        \begin{tabular}{@{} l @{}}
            \textbf{Inversión Angular (PDoA):} \\[3pt]
            Retardo espacial: \ $\Delta r = d \sin\theta$ \\
            Desfase medido: \ $\Delta\Phi = \frac{2\pi f_0}{c(T)} d \sin\theta$ \\[5pt]
            \textcolor{UdeAGreen!80!black}{\textbf{Cota de no-aliasing espacial:}} \\[2pt]
            $d \le \frac{\lambda_0}{2} = \frac{c(T)}{2 f_0}$
        \end{tabular}
    };

\end{tikzpicture}